\section{Introduction}\label{sec:intro}

Suppose you, the player, are shown a grid of points with $h$ rows and $m$ columns. Each
point is either white or blacked out, and a blacked-out point cannot tell you anything.
Before the game starts you choose which points to black out. We call the set of
blacked-out points the \emph{blackout}.

Then the game master draws a rectangle with its four corners on grid points. You cannot
see the rectangle. You are only told which of its corners are white. The rectangle does
not have to line up with the grid. Tilted rectangles count too. You must be able to say
where all four corners are with complete confidence, whatever rectangle the game master
draws. Kagey's question is how many points you can black out and still do this. The order
of play matters. The blackout is fixed before the rectangle is drawn, so a blackout is
only good if \emph{no} two rectangles could ever look the same to you.

We write $(i,j)$ for the point in column $i$ and row $j$, counting columns from the left
and rows from the top. Figure~\ref{fig:game} plays one round on the $3\times4$ grid with
row $1$ and column $1$ blacked out. The game master draws the tilted square with corners
$(2,2)$, $(3,1)$, $(4,2)$, and $(3,3)$. We are shown $(2,2)$, $(4,2)$, and $(3,3)$.
Three corners of a rectangle determine the fourth (Lemma~\ref{lem:threecorners}), so the
fourth corner is $(3,1)$. This blackout of six points works against every rectangle, and
no good blackout of the $3\times4$ grid has more points (Section~\ref{sec:three}).

\begin{figure}[t]
\centering
\resizebox{0.7\textwidth}{!}{\begin{tikzpicture}[x=0.9cm,y=0.9cm,baseline=(current bounding box.north)]
  \begin{scope}[xshift=0.000cm]
    \draw[black,line width=0.36pt,dash pattern=on 1.35pt off 1.35pt] (1,-1) -- (4,-1) (1,-2) -- (4,-2) (1,-3) -- (4,-3) (1,-1) -- (1,-3) (2,-1) -- (2,-3) (3,-1) -- (3,-3) (4,-1) -- (4,-3);
    \filldraw[fill=white,draw=black,line width=0.54pt] (1,-1) circle (0.3);
    \filldraw[fill=white,draw=black,line width=0.54pt] (2,-1) circle (0.3);
    \filldraw[fill=white,draw=black,line width=0.54pt] (3,-1) circle (0.3);
    \filldraw[fill=white,draw=black,line width=0.54pt] (4,-1) circle (0.3);
    \filldraw[fill=white,draw=black,line width=0.54pt] (1,-2) circle (0.3);
    \filldraw[fill=white,draw=black,line width=0.54pt] (2,-2) circle (0.3);
    \filldraw[fill=kyellow,draw=black,line width=0.54pt] (3,-2) circle (0.3);
    \filldraw[fill=white,draw=black,line width=0.54pt] (4,-2) circle (0.3);
    \filldraw[fill=white,draw=black,line width=0.54pt] (1,-3) circle (0.3);
    \filldraw[fill=white,draw=black,line width=0.54pt] (2,-3) circle (0.3);
    \filldraw[fill=white,draw=black,line width=0.54pt] (3,-3) circle (0.3);
    \filldraw[fill=white,draw=black,line width=0.54pt] (4,-3) circle (0.3);
    \node[lab] at (1,-0.45) {1};
    \node[lab] at (2,-0.45) {2};
    \node[lab] at (3,-0.45) {3};
    \node[lab] at (4,-0.45) {4};
    \node[lab] at (0.45,-1) {1};
    \node[lab] at (0.45,-2) {2};
    \node[lab] at (0.45,-3) {3};
    \node[capt,anchor=north] at (2.50,-3.45) {(a) the point $(3,2)$};
  \end{scope}
  \begin{scope}[xshift=4.680cm]
    \draw[black,line width=0.36pt,dash pattern=on 1.35pt off 1.35pt] (1,-1) -- (4,-1) (1,-2) -- (4,-2) (1,-3) -- (4,-3) (1,-1) -- (1,-3) (2,-1) -- (2,-3) (3,-1) -- (3,-3) (4,-1) -- (4,-3);
    \filldraw[fill=kgray,draw=black,line width=0.54pt] (1,-1) circle (0.3);
    \draw[black,line width=0.54pt] (0.788,-0.788) -- (1.212,-1.212) (0.788,-1.212) -- (1.212,-0.788);
    \filldraw[fill=kgray,draw=black,line width=0.54pt] (2,-1) circle (0.3);
    \draw[black,line width=0.54pt] (1.788,-0.788) -- (2.212,-1.212) (1.788,-1.212) -- (2.212,-0.788);
    \filldraw[fill=kgray,draw=black,line width=0.54pt] (3,-1) circle (0.3);
    \draw[black,line width=0.54pt] (2.788,-0.788) -- (3.212,-1.212) (2.788,-1.212) -- (3.212,-0.788);
    \filldraw[fill=kgray,draw=black,line width=0.54pt] (4,-1) circle (0.3);
    \draw[black,line width=0.54pt] (3.788,-0.788) -- (4.212,-1.212) (3.788,-1.212) -- (4.212,-0.788);
    \filldraw[fill=kgray,draw=black,line width=0.54pt] (1,-2) circle (0.3);
    \draw[black,line width=0.54pt] (0.788,-1.788) -- (1.212,-2.212) (0.788,-2.212) -- (1.212,-1.788);
    \filldraw[fill=white,draw=black,line width=0.54pt] (2,-2) circle (0.3);
    \filldraw[fill=white,draw=black,line width=0.54pt] (3,-2) circle (0.3);
    \filldraw[fill=white,draw=black,line width=0.54pt] (4,-2) circle (0.3);
    \filldraw[fill=kgray,draw=black,line width=0.54pt] (1,-3) circle (0.3);
    \draw[black,line width=0.54pt] (0.788,-2.788) -- (1.212,-3.212) (0.788,-3.212) -- (1.212,-2.788);
    \filldraw[fill=white,draw=black,line width=0.54pt] (2,-3) circle (0.3);
    \filldraw[fill=white,draw=black,line width=0.54pt] (3,-3) circle (0.3);
    \filldraw[fill=white,draw=black,line width=0.54pt] (4,-3) circle (0.3);
    \node[lab] at (1,-0.45) {1};
    \node[lab] at (2,-0.45) {2};
    \node[lab] at (3,-0.45) {3};
    \node[lab] at (4,-0.45) {4};
    \node[lab] at (0.45,-1) {1};
    \node[lab] at (0.45,-2) {2};
    \node[lab] at (0.45,-3) {3};
    \node[capt,anchor=north] at (2.50,-3.45) {(b) our blackout};
  \end{scope}
  \begin{scope}[xshift=9.360cm]
    \draw[black,line width=0.36pt,dash pattern=on 1.35pt off 1.35pt] (1,-1) -- (4,-1) (1,-2) -- (4,-2) (1,-3) -- (4,-3) (1,-1) -- (1,-3) (2,-1) -- (2,-3) (3,-1) -- (3,-3) (4,-1) -- (4,-3);
    \draw[kred,line width=2.55pt,line join=round] (2,-2) -- (3,-1) -- (4,-2) -- (3,-3) -- cycle;
    \filldraw[fill=kgray,draw=black,line width=0.54pt] (1,-1) circle (0.3);
    \draw[black,line width=0.54pt] (0.788,-0.788) -- (1.212,-1.212) (0.788,-1.212) -- (1.212,-0.788);
    \filldraw[fill=kgray,draw=black,line width=0.54pt] (2,-1) circle (0.3);
    \draw[black,line width=0.54pt] (1.788,-0.788) -- (2.212,-1.212) (1.788,-1.212) -- (2.212,-0.788);
    \filldraw[fill=kgray,draw=black,line width=0.54pt] (3,-1) circle (0.3);
    \draw[black,line width=0.54pt] (2.788,-0.788) -- (3.212,-1.212) (2.788,-1.212) -- (3.212,-0.788);
    \filldraw[fill=kgray,draw=black,line width=0.54pt] (4,-1) circle (0.3);
    \draw[black,line width=0.54pt] (3.788,-0.788) -- (4.212,-1.212) (3.788,-1.212) -- (4.212,-0.788);
    \filldraw[fill=kgray,draw=black,line width=0.54pt] (1,-2) circle (0.3);
    \draw[black,line width=0.54pt] (0.788,-1.788) -- (1.212,-2.212) (0.788,-2.212) -- (1.212,-1.788);
    \filldraw[fill=white,draw=black,line width=0.54pt] (2,-2) circle (0.3);
    \filldraw[fill=white,draw=black,line width=0.54pt] (3,-2) circle (0.3);
    \filldraw[fill=white,draw=black,line width=0.54pt] (4,-2) circle (0.3);
    \filldraw[fill=kgray,draw=black,line width=0.54pt] (1,-3) circle (0.3);
    \draw[black,line width=0.54pt] (0.788,-2.788) -- (1.212,-3.212) (0.788,-3.212) -- (1.212,-2.788);
    \filldraw[fill=white,draw=black,line width=0.54pt] (2,-3) circle (0.3);
    \filldraw[fill=white,draw=black,line width=0.54pt] (3,-3) circle (0.3);
    \filldraw[fill=white,draw=black,line width=0.54pt] (4,-3) circle (0.3);
    \node[lab] at (1,-0.45) {1};
    \node[lab] at (2,-0.45) {2};
    \node[lab] at (3,-0.45) {3};
    \node[lab] at (4,-0.45) {4};
    \node[lab] at (0.45,-1) {1};
    \node[lab] at (0.45,-2) {2};
    \node[lab] at (0.45,-3) {3};
    \node[capt,anchor=north] at (2.50,-3.45) {(c) the rectangle};
  \end{scope}
  \begin{scope}[xshift=14.040cm]
    \draw[black,line width=0.36pt,dash pattern=on 1.35pt off 1.35pt] (1,-1) -- (4,-1) (1,-2) -- (4,-2) (1,-3) -- (4,-3) (1,-1) -- (1,-3) (2,-1) -- (2,-3) (3,-1) -- (3,-3) (4,-1) -- (4,-3);
    \filldraw[fill=kyellow,draw=black,line width=0.54pt] (2,-2) circle (0.3);
    \filldraw[fill=kyellow,draw=black,line width=0.54pt] (4,-2) circle (0.3);
    \filldraw[fill=kyellow,draw=black,line width=0.54pt] (3,-3) circle (0.3);
    \node[lab] at (1,-0.45) {1};
    \node[lab] at (2,-0.45) {2};
    \node[lab] at (3,-0.45) {3};
    \node[lab] at (4,-0.45) {4};
    \node[lab] at (0.45,-1) {1};
    \node[lab] at (0.45,-2) {2};
    \node[lab] at (0.45,-3) {3};
    \node[capt,anchor=north] at (2.50,-3.45) {(d) what we are shown};
  \end{scope}
\end{tikzpicture}
}
\caption{One round on the $3\times4$ grid. Black dots are blacked out and open dots are
white. In (d) the gray specks are points we are told nothing about.}
\label{fig:game}
\end{figure}

Kagey poses the problem in his Open Problem Collection~\cite{kagey001}. It is also
Problem 1 of the collection he handed out at the fourteenth Gathering 4
Gardner~\cite{kageyG4G14}. We write $f(h,m)$ for the largest size of a blackout of the
$h\times m$ grid that keeps every rectangle identifiable. The number of such blackouts of
that size, counted with their coordinates, is $a_h(m)$. These are the tables
\seqnum{A397315} and \seqnum{A399580} of the OEIS~\cite{oeis}, which we contributed. Our
comments on \seqnum{A397315} record $f(2,m)=m+1$, the conjecture $f(3,m)=m+2$, and the
guess that $f(h,m)=h+m$ whenever $\min(h,m)\ge4$, apart from $f(4,4)=7$ and $f(4,5)=8$.
The guess was based on seven computed cells, and we show here that it fails. The
three-row conjecture was proved in Lean by The Omega Institute~\cite{trureturing2026}, by
a finite-set inclusion-exclusion argument, on 15 September 2026. It is also the case
$h=3$ of our row bound (Theorem~\ref{thm:rowbound}), and we do not claim it as new. We
will update both entries.

\begin{example}[the $2\times3$ strip]\label{ex:strip23}
On two rows every rectangle is axis-aligned, and $R_{a,b}$ is the one on columns $a<b$.
Black out $(3,1)$, $(1,2)$, $(2,2)$, and $(3,2)$. Then $R_{1,2}$ shows $(1,1)$ and
$(2,1)$, $R_{1,3}$ shows $(1,1)$, and $R_{2,3}$ shows $(2,1)$, so this blackout is good.
If we black out columns $1$ and $2$ instead, then $R_{1,3}$ and $R_{2,3}$ both show
$(3,1)$ and $(3,2)$. With five black points every view is empty or the one kept point, so
two of the three rectangles look the same. Hence $f(2,3)=4$, and $a_2(3)=12$
(Theorem~\ref{thm:strip}).
\end{example}

Two rectangles look the same exactly when every corner that belongs to only one of them
is black (Lemma~\ref{lem:diff}). So everything turns on the rectangles that show at most
two corners.

\paragraph{Results.} Put $q_h=\lfloor h^2/4\rfloor$. Computer checks are marked in each
statement.

\begin{thmA}[narrow grids]
\begin{enumerate}[(a)]
\item For all $h,m\ge1$ we have $f(h,m)\ge h+m-1$, and for $m\ge3$ we have
$f(h,m)\le m+q_h$.
\item For $m\ge3$ we have $f(2,m)=m+1$ and $f(3,m)=m+2$. Also $f(4,m)=m+4$ for every
$m\ge6$, and $f(5,5)=10$.
\item (Computer-assisted.) We have $f(4,4)=7$, $f(4,5)=8$, $f(5,6)=11$, and $f(5,7)=12$, while
$f(5,m)=m+6$ for every $m\ge8$.
\item (Computer-assisted for $m\ge9$.) For $8\le m\le14$ we have $f(6,m)\le m+8$, and
$f(6,m)=m+9$ for every $m\ge15$.

\item (Computer-assisted.) We have $f(6,6)=12$, $f(6,7)=13$, $f(7,7)=14$, and $f(6,8)=15$.
\end{enumerate}
\end{thmA}

The lower bound in (a) is the boundary cross of Lemma~\ref{lem:cross}. The upper bound is
the row bound of Theorem~\ref{thm:rowbound}, and its proof rests on Mantel's theorem.
Parts (b) to (e) collect Theorems~\ref{thm:strip}, \ref{thm:threevalue},
\ref{thm:fourvalue}, and~\ref{thm:fivefive}, Corollary~\ref{cor:fourboundary}, and
Section~\ref{sec:five}. So for $h\le6$ the row bound is attained at every large width.
The guess $f(h,m)=h+m$ first fails on the $5\times8$ grid (Theorem~\ref{thm:refute}).

\begin{thmB}[three-row optima]
For $m\ge4$ we have $a_3(m)=4(m-1)(3m+2)\,3^{m-5}+E_m$, where $E_m\ge0$ counts the optima
with a column black in rows $1$ and $3$. So $a_3(m)\sim\frac4{81}m^23^m$, and a fraction
tending to $4/27$ of the $m^23^{m-1}$ axis-valid candidates is valid. The series
$\sum_{m\ge3}a_3(m)z^m$ is rational. (Computer-assisted.) Its reduced denominator is
$(1-3z)^3(1-2z-2z^2)^3(1+2z^2)^3$, and $E_m$ has a closed form with
$E_m\sim\frac{19+11\sqrt3}{432}m^2(1+\sqrt3)^m$.
\end{thmB}

So almost all three-row optima have two doubles in the rows $\{1,2\}$ and $\{2,3\}$, and
the others grow only like $(1+\sqrt3)^m$ (Theorems~\ref{thm:threecount}
and~\ref{thm:threeclosed}). The series is Theorem~\ref{thm:threegf}, with the proof
finished in Appendix~\ref{app:threegf}, and the asymptotics are
Corollary~\ref{cor:threeasymp}. The sequence starts $53,224,892,3420,12704$ at $m=3$.

\begin{thmC}[four rows]
Let $m\ge4$. No valid blackout of the $4\times m$ grid hides all four corners of a
rectangle. For $m\ge6$,
\[
a_4(m)\le3\,(m)_4\,4^{m-4}\Bigl(\frac{53}{64}\Bigr)^{\max(0,\lfloor m/4\rfloor-4)},
\qquad (m)_4=m(m-1)(m-2)(m-3).
\]
(Computer-assisted.) The series $\sum_{m\ge6}a_4(m)z^m$ is rational with reduced
denominator of degree $420$, and $a_4(m)\sim c\,m^4\rho_4^m$ with $c=8.1611\times10^{-5}$
and $\rho_4=3.35199\ldots$, an algebraic number of degree $20$.
\end{thmC}

So valid blackouts become rare on four rows, since the bound has the base
$212^{1/4}=3.8157\ldots<4$ (Theorems~\ref{thm:hidden} and~\ref{thm:foursurvival}). The
generating function and the growth come from a transfer matrix with a few thousand states
(Theorem~\ref{thm:fourgfexact} and Corollary~\ref{cor:fourgrowth}).


\begin{thmD}[finite collar]
Let $h\ge3$ and $W_h=1+\lfloor(h-1)^2/4\rfloor$. Pad a blackout of an $h\times m$ grid, on
one side or on both, by columns black only in the top row, or only in the bottom row.
Then every such padding is valid if and only if the padding by $W_h$ columns on each
chosen side is valid. (Computer-assisted.) A blackout of the $5\times8$ grid and one of the
$6\times16$ grid pass this test with the top row on the right. Hence $f(5,m)=m+6$ for
$m\ge8$ and $f(6,m)=m+9$ for $m\ge16$, and a stored blackout of the $6\times15$ grid gives
$f(6,15)=24$.
\end{thmD}

The collar theorem is Theorem~\ref{thm:collar}, and Section~\ref{sec:five} has the two
blackouts, with test boards of widths $13$ and $23$. The padding row has to be a boundary
row (Remark~\ref{rem:interior}).

\begin{thmE}[every fixed height]
Fix $h\ge1$ and write a blackout of the $h\times m$ grid as the word $B_1\cdots B_m$,
where $B_i$ is the set of black rows of column $i$. The valid words form a regular
language. If $h\ge2$, the excess $f(h,m)-m$ is eventually periodic in $m$, and
$\sum_{m\ge0}a_h(m)z^m$ is rational.
\end{thmE}

These are Theorems~\ref{thm:regular} and~\ref{thm:periodic}. They give no period and no
threshold. Note that validity is not a property of bounded windows
(Proposition~\ref{prop:nonlocal}), so the proof reads two rectangles in parallel with the
word.

\paragraph{Prior work.} In its hitting-set form (Corollary~\ref{cor:hitting}) the game
asks for a test cover in the sense of Katona~\cite{katona1966}, Bellitto~\cite{bellitto2018},
and Foucaud~\cite[Section~1.1]{foucaud2015}. The axis-aligned rectangles alone lead to
bipartite graphs without cycles of length four and six (Lemma~\ref{lem:girth}).
Erd\H{o}s~\cite{erdos1080} and de Caen and Sz\'ekely~\cite{decaenszekely1997} study such
graphs. Neuwirth~\cite[Theorem~1 and Proposition~17]{neuwirth2001} proves our row bound
for them when $m\ge q_h$, with the extremal graphs of Theorem~\ref{thm:equality} as
examples. Our proof uses Mantel's theorem~\cite{mantel1907}, covers every $m\ge3$, and
characterizes the equality case with its count. None of these works
sees tilted rectangles.

Euler~\cite{euler2005}, Butler, Choi, Kim, and Seo~\cite{butler2019},
Mertens~\cite{mertens2024}, and Alanko, Crevals, Isopoussu, \"Osterg\r{a}rd, and
Pettersson~\cite{alanko2011} treat grids of fixed height by column
transfer~\cite[Proposition~V.6]{flajolet2009}. Jiang~\cite{jiang2018} proves that optimal
identifying codes in strips can be taken periodic. Theorem~E is of the same kind.
Guilhot, Little, and Parkinson~\cite[Theorem~1]{glp2025} prove that the maximizers of a
bounded weight on a regular language form a regular language. Baccelli, Cohen, Olsder,
and Quadrat~\cite{baccelli1992} prove that max-plus matrix powers are eventually
periodic. We maximize separately at each width, so we prove Theorem~E directly
(Section~\ref{sec:height}).

Section~\ref{sec:conv} fixes the conventions and proves the difference lemma.
Section~\ref{sec:strip} solves the strips, and Section~\ref{sec:graph} proves the general
bounds and lists the tilted rectangles of narrow grids. Sections~\ref{sec:three}
to~\ref{sec:five} treat three, four, five, and six rows, and Section~\ref{sec:height}
treats every fixed height. Section~\ref{sec:verify} describes the computations, code, and
data~\cite{repo}, and Section~\ref{sec:open} lists open questions.
